\section{总结}
当格点规范理论中的抽样发生临界慢化时，在理论的临界极限附近精确估计人们感兴趣的物理量就会受到阻碍。基于流的模型使我们能够直接从目标分布的一个近似中抽样，由此导出的物理可观测量估计值在无限统计极限下是 exact 的。本文引入了构造上满足精确规范不变性的基于流的模型，并发现把该方法应用于二维阿贝尔规范理论时，对拓扑量的估计比 HMC、热浴等现有算法更为高效。

这里提出的方法一般地适用于由李群定义的规范理论，包括 QCD 这类非阿贝尔理论。要把该方法推广到这类理论，必须把有表达力的可逆函数定义为规范等变耦合层的核。有几条可能的前进路径；例如文献~\cite{rezende2020normalizing} 定义了球面流形上的流，文献~\cite{falorsi19lie} 定义了李群上的满射（虽然不是双射）映射，二者都使用了神经网络参数化。未来的工作将探索在这些方法的推广之上构造核，从而为 QCD 这样的非阿贝尔理论产生规范不变的流。
